\documentclass[12pt,a4paper]{article}
\usepackage[utf8]{inputenc}
\usepackage[T1]{fontenc}
\usepackage{lmodern}

\usepackage[margin=2.5cm]{geometry}
\usepackage{graphicx}
\usepackage{float}
\usepackage{booktabs}
\usepackage{longtable}
\usepackage{enumitem}
\usepackage{xcolor}
\usepackage{hyperref}
\usepackage{listings}

\definecolor{codegray}{gray}{0.95}

\lstset{
  backgroundcolor=\color{codegray},
  basicstyle=\ttfamily\small,
  breaklines=true,
  frame=single
}

\hypersetup{
  colorlinks=true,
  linkcolor=blue,
  urlcolor=blue,
  citecolor=blue
}

% ------------------------------------------------------------
% Header del documento
% ------------------------------------------------------------
\title{
  \textbf{Sistema de autenticación de usuarios}\\
  \large Documento del diseño de la arquitectura
}

\author{
  Juan Felipe Rodríguez Amador
}

\date{\today}

\begin{document}

\maketitle

% --------
% Tabla de contenidos
% --------

\tableofcontents
\newpage

% ============================================================
\section{Descripción del Sistema}
% ============================================================

La sistema busca autenticar usuarios priorizando una alta
disponibilidad y un despliegue sencillo para el usuario. Esto
mediante una aplicación móvil,
haciendo uso de Spring Boot como el backend que se conecta a una base
de datos de MySQL.

% ============================================================
\section{Atributos de Calidad}
% ============================================================

La arquitectura prioriza los siguientes atributos de calidad:

\begin{itemize}
  \item Alta disponibilidad (de un 99.99\%)
  \item Despliegue sencillo
\end{itemize}

% ============================================================
\section{Descripción general de la Arquitectura}
% ============================================================

El sistema sigue una arquitectura de 3-niveles compuesta por un nivel
ininicial de presentación, luego uno de lógica de negocio y por
último uno de datos.

\subsection{Diseño de la Arquitectura}

\begin{figure}[H]
  \centering
  \includegraphics[width=0.9\textwidth]{./architecture.png}
  \caption{Diagrama lógico de la arquitectura}
  \label{fig:architecture}
\end{figure}

% ============================================================
\section{Niveles de la Arquitectura}
% ============================================================

\subsection{Nivel de presentación}

Se implementa una aplicación móvil de android nativo desarrollada en
el framework de Jetpack Compose...

\subsection{Nivel de logica de negocio}

Se hace uso de un backend desarrollado en Spring Boot...

\subsection{Nivel de datos}

Se hace uso de una base de datos de MySQL

% ============================================================
\section{Tacticas Utilizadas}
% ============================================================

Se utilizaron las siguientes tácticas para cumplir con los atributos
de calidad...

\subsection{Alta disponibilidad}

Para garantizar una alta disponivlidad se utilizó...

\begin{itemize}
  \item Redundancia
  \item Detección de fallos
  \item Ping/echo
\end{itemize}

\subsection{Despliegue sencillo}

Para garantizar un despliegue sencillo se utilizó...

\begin{itemize}
  \item Contenedores
\end{itemize}

\end{document}
